\chapter{From a circular torus to a min-max theorem}\label{sec:history-willmore-development}
The sources here are the authors' papers: \cite{LY82}, the opening pages of \cite{SimW93}, and the introduction of \cite{MNWill}. The entries distinguish the historical questions of identifying a candidate, proving existence of an optimizer, and determining the sharp value. They do not replace the mathematical statements and proof dependencies linked below.

\begin{remark}[A family calculation suggested a global conjecture]\label{rem:history-willmore-circular-candidate}
Marques--Neves date Willmore's torus conjecture to 1965 in \cite{MNWill}, Section 1.1, p.~684. The circular torus with ratio $R/r=\sqrt2$ attains energy $2\pi^2$, but a calculation over circular tori leaves all other immersions untreated. The same candidate is a stereographic image of the Clifford torus in $S^3$. Theorems \ref{thm:circular-ring-willmore-energy} and \ref{thm:willmore-first-variation-spaceform}, and Example \ref{ex:clifford-stereographic-ring}, explain why this candidate is natural without assuming the eventual global theorem.
\end{remark}

\begin{remark}[Conformal area separated the multiple-point obstruction]\label{rem:history-li-yau-multiplicity}
In their 1982 paper \cite{LY82}, Li--Yau introduced conformal volume and proved a multiplicity bound for the bending energy of an immersed surface; see the introduction and Theorem 6, p.~286. A double point forces energy at least $8\pi$, already above the conjectured torus minimum $2\pi^2$. Thus the unresolved low-energy problem concerned embeddings. Theorem \ref{thm:li-yau-smooth-multiplicity} gives a direct radial proof of that obstruction, while Corollary \ref{cor:immersed-positive-genus-willmore} shows precisely how it combines with the later embedded theorem. The two numerical thresholds serve different purposes.
\end{remark}

\begin{remark}[Existence and the exact minimum were separate achievements]\label{rem:history-willmore-existence-value}
Simon's 1993 paper \cite{SimW93}, pp.~281--283, established a smooth embedded torus minimizing Willmore energy among embedded tori in each Euclidean ambient dimension at least three. This supplied an optimizer without identifying its sharp energy in three-space. The 2014 paper of Marques--Neves \cite{MNWill}, Theorem A, p.~685, determined the positive-genus bound in $S^3$ and its equality case through min-max theory. Their Corollary D, p.~687, realizes the Clifford torus through a five-parameter min-max family. The local square identities developed here illuminate Simon's estimates; neither his full existence proof nor the five-parameter min-max proof has been replaced by those local calculations.
\end{remark}
